\subsection{The exact ramification threshold for a mixed cyclic comparison}

\begin{lemma}[A rational cyclic point at arbitrary level]
\label{lem:B-RD63-point}
\status{Candidate; complete Newton-polygon proof written, independent audits pending.}
Let $p\ge5$, let $R$ be a complete discrete valuation ring of mixed
characteristic $(0,p)$ with algebraically closed residue field, and
normalize $v(\pi_R)=1$. Put $a=v(p)$ and assume $p\nmid a$.
If a one-dimensional formal group over $R$ of special height two
has an $R$-rational point of exact order $p^n$, where $n\ge2$, then
\[
                    a>(p^2-1)p^{2n-3}.                 \tag{B-RD63.1}
\]
\end{lemma}
\begin{proof}
Use the integral universal coordinate and coefficient bounds in
the full proof of Lemma~\ref{lem:B-RD49-point}; these come from
\cite[Proposition~5.7; (12.3)--(12.5); Proposition~12.10 and proof]
{HopkinsGross1994Bundles}. Put $q=p^2$, $N=q-1$, and $b=v(u)$,
allowing $b=+\infty$. The only possible vertices of the positive
root Newton polygon of $[p](T)$ are
\[
                         (1,a),\ (p,b),\ (q,0).
\]
All intermediate coefficients lie in $(p,u)$, coefficients in
degrees $2,\ldots,p-1$ vanish, and the degree-$q$ coefficient is a
unit. Later integral coefficients lie strictly above the relevant
negative-slope segments. Thus, if $b\ge ap/(p+1)$, every nonzero
$p$-torsion point has valuation $a/N$. Otherwise its two valuations
are
\[
       \alpha=\frac{a-b}{p-1},\qquad
       \beta=\frac{b}{p(p-1)},                           \tag{B-RD63.2}
\]
with multiplicities $p-1$ and $q-p$, respectively.

We also require the exact inverse-step polygon. For any nonzero
torsion point of valuation $Z<a$, the polynomial obtained by
Weierstrass preparation from $[p](T)-z$ has its possible lower
vertices at $(0,Z),(p,b),(q,0)$; the point $(1,a)$ is strictly
above them. Consequently its $p$-preimages have valuations
\[
 \begin{cases}
 Z/q &\text{if } b\ge Z(p-1)/p,\\
 (Z-b)/p\ \text{($p$ times)},\quad
 \beta\ \text{($q-p$ times)}&\text{if }b<Z(p-1)/p.
 \end{cases}                                           \tag{B-RD63.3}
\]
At equality both displayed slopes agree. Every inverse-step
valuation is strictly smaller than $Z$. Every nonzero $p$-torsion
valuation is at most $a/(p-1)<a$, so (B-RD63.3) applies at every
step of a rational chain. Once a valuation is $\beta$, all further
inverse steps divide it by $q$, because
$b=p(p-1)\beta>\beta(p-1)/p$. The same remains true after any
uniform inverse step whose input satisfies the first inequality
of (B-RD63.3).

Let $P$ have order $p^n$ and let $Z$ be the valuation of
$[p]^{n-1}P$. All points in this chain are rational, so their
positive valuations are integers. In the one-slope case all
inverse steps divide by $q$, giving
$v(P)=a/(Nq^{n-1})$. Its integrality contradicts $p\nmid a$.
In the two-slope case, if $Z=\beta$, then
\[
 b\ge p(p-1)q^{n-1},\qquad
 a>\frac{p+1}{p}b\ge Nq^{n-1},
\]
which is stronger than (B-RD63.1).

It remains that $Z=\alpha$. If $b\ge\alpha(p-1)/p$, all inverse
steps are uniform, and integrality gives $\alpha\ge q^{n-1}$.
Hence
\[
 a=b+(p-1)\alpha\ge\frac{N}{p}\alpha
       \ge Np^{2n-3}.
\]
Equality would make $p\mid a$, so the desired inequality is strict.
If instead $b<\alpha(p-1)/p$, a rational high inverse step of
valuation $A=(\alpha-b)/p$ would imply
\[
                   a=p\bigl(b+(p-1)A\bigr),
\]
again impossible. The first inverse step must therefore have
valuation $\beta$. The remaining $n-2$ steps divide by $q$;
thus $b\ge p(p-1)q^{n-2}$. The strict inequality for this case
is equivalent to $a>(p+1)b$, and gives
$a>Np^{2n-3}$. This exhausts both Newton polygons and every
rational inverse-step choice.
\end{proof}

\begin{proposition}[The mixed-case necessary threshold]
\label{prop:B-RD63-mixed}
\status{Candidate; complete reduction written, independent audits pending.}
Let $K$ be complete discretely valued of mixed characteristic
$(0,p)$, $p\ge5$, with perfect residue field and absolute
ramification $e$ prime to $p$. Let $M$ be a height-two BT group
which acquires one-dimensional supersingular good reduction,
and let $H$ admit a filtration by two potentially good height-one
groups whose dimensions sum to one. For $n\ge2$,
\[
              M[p^n]\simeq H[p^n]\quad\Longrightarrow
                         e>p^{2n-3}.                   \tag{B-RD63.4}
\]
\end{proposition}
\begin{proof}
Pass to the completed maximal unramified extension, preserving
$e$ and the torsion comparison. Its residue field is algebraically
closed. The full quotient-and-twist proof of
Proposition~\ref{prop:B-RD49-mixed} works at level $p^n$ without
change: transport the first filtered subgroup to a primitive
cyclic subgroup of $M[p^n]$. If its character is a finite
$\epsilon$, twist by $\epsilon^{-1}$. If its character is
$\chi_p\epsilon$, first quotient by that cyclic subgroup and
then twist by the inverse finite quotient character. This produces
a potentially supersingular height-two group with a rational
point of exact order $p^n$. Finite-flat closure in a good model
justifies the quotient, and the determinant is $\chi_p$ times
a finite $\mu_{p-1}$-valued character, exactly as in that proof.

Lemma~\ref{lem:B-RD49-inertia} gives good formal reduction over
a tame extension of degree $m\mid p^2-1$. Its absolute index is
$a=me$, prime to $p$. The rational point extends to the finite
flat torsion model and lies in its formal group, since the special
fibre is connected. Lemma~\ref{lem:B-RD63-point} now gives
$me>(p^2-1)p^{2n-3}$. Since $m\le p^2-1$, (B-RD63.4) follows.
\end{proof}

\begin{theorem}[All indices above the cyclic threshold occur]
\label{thm:B-RD63-arc}
\status{Candidate; complete extension of the audited arc construction written.}
Let $p\ge5$, $n\ge2$, and let $e$ be a positive multiple of
$p-1$ with $e>p^{2n-3}$. There is a complete mixed-characteristic
field $K$ with residue field $\overline{\mathbf F}_p$, absolute
ramification $e$, and $\mu_p\subset K$, and height-two BT groups
$M,H$ such that $M[p^n]\simeq H[p^n]$, while their requested
prime-to-$p$ special characters differ. Here $M$ is wholly
potentially supersingular, and $H$ has a two-step potentially
good height-one filtration of total dimension one.
\end{theorem}
\begin{proof}
Put $q=p^2$, $N=q-1$, $Q=q^{n-1}$, and $x_0=e/Q$.
Choose an integer $A$ with
\[
                  x_0<A<px_0,\qquad N\nmid A.          \tag{B-RD63.5}
\]
For completeness take $A_0=\lfloor x_0\rfloor+1$, and replace it
by $A_0+1$ only if $N\mid A_0$. If $x_0<1$, the hypothesis
$px_0>1$ makes $A_0=1$ valid. If $x_0\ge1$, the chosen value
satisfies $A\le x_0+2<px_0$ because $p\ge5$; two consecutive
integers cannot both be divisible by $N$.

Use the universal group and exact logarithm from the proof of
Lemma~\ref{lem:B-RD62-arc}, and define
\[
 \begin{gathered}
 d=Ne,\qquad b=d-(p-1)AQ,\qquad Y_0(X)=X^A,\\
 Y_{j+1}(X)=[p]_{X^b}(Y_j(X)),\qquad
 P_A(X)=\Phi_{X^b}(Y_{n-1}(X)).
 \end{gathered}                                        \tag{B-RD63.6}
\]
The two strict inequalities in (B-RD63.5) give exactly
\[
 \begin{gathered}
 b>(q-p)Aq^{n-2}>0,\qquad AQ<d<NAQ,\\
 b\equiv-(p-1)A\pmod N.
 \end{gathered}                                        \tag{B-RD63.7}
\]
Indeed the first difference is $N(e-Ap^{2n-3})$; also
$AQ<pe<Ne=d$. Every coefficient argument of the audited arc
now applies with all input weights multiplied by $A$:
\[
 Y_j(X)\equiv X^{Aq^j}+\text{higher terms}\pmod p,
 \quad P_A(0)=p,
 \quad P_A(X)\equiv X^d+\text{higher terms}\pmod p.
\]
Thus $R=W(\overline{\mathbf F}_p)[[X]]/(P_A)$ is a complete
Eisenstein DVR with uniformizer $x$ and $v_R(p)=d$.
The strict bound $d>AQ$ prevents $p$-divisible coefficients from
changing any valuation $v_R(Y_j(x))=Aq^j$. Consequently $x^A$
has exact order $p^n$ in $F_{x^b}$.

Choose $\zeta$ of order $N$ and set $a=\zeta^{-A}$.
The exact homogeneous logarithm identities (B-RD62.9)--(B-RD62.10)
hold for this Teichmuller unit too. The substitution
$\sigma(x)=\zeta x$ sends $u=x^b$ to $a^{p-1}u$, sends every
$Y_j$ to $a^{-1}Y_j$, and fixes $P_A$. It has order $N$ on $R$.
The semilinear coordinate map $aT$ fixes $x^A$ and satisfies
the descent cocycle. Therefore the generic BT group and marked
point descend to $K=\operatorname{Frac}(R)^{\langle\sigma\rangle}$,
with $e_K=e$. This degree-$N$ good extension need not be faithful
on the special group; no faithfulness is required here.
The same tame cyclotomic argument as in B-RD62 gives $\mu_p\subset K$.

On the special isocrystal the two eigenvalues are $a,a^p$.
The determinant is $\chi_p\psi$, where
$\psi(\sigma)=a^{p+1}$ has order dividing $p-1$.
Apply the identical finite-flat cyclic quotient, twist by
$\psi^{-1}$, and compatible Kummer-extension lifting of
Theorem~\ref{thm:B-RD62} and Lemma~\ref{lem:B-RD62-extension}.
They produce $M,H$ with exactly the claimed common truncation.
Writing $\lambda=\zeta^A$, their special eigenvalue multisets are
\[
                       \{\lambda,\lambda^p\},\qquad
                       \{1,\lambda^{p+1}\}.
\]
Their trace difference is $-(\lambda-1)(\lambda^p-1)\ne0$, since
$N\nmid A$ and $p$ is coprime to $N$. This is a difference on
an element of prime-to-$p$ order in the finite special inertia
quotient, with no assertion of finite order in absolute inertia.
\end{proof}

\begin{corollary}[Sharp mixed threshold and a general lower bound]
\label{cor:B-RD63}
\status{Candidate; exact scope and indexing written, independent audits pending.}
For fixed $p\ge5$, $n\ge2$, and a positive multiple $e$ of $p-1$
prime to $p$, a comparison of the two classes in
Proposition~\ref{prop:B-RD63-mixed} can occur over some permitted
field of index $e$ if and only if $e>p^{2n-3}$.
For every positive multiple $e$ of $p-1$ and height $h\ge2$,
the universal B.4 index satisfies
\[
 \nu_{\min}(p,e,h)\ge
 \left\lfloor\frac{\lfloor\log_p(e-1)\rfloor+3}{2}\right\rfloor.
                                                               \tag{B-RD63.8}
\]
\end{corollary}
\begin{proof}
Necessity and existence are exactly the preceding proposition
and theorem; the latter is an existence assertion for some $K$,
not for every field with that index. A failure at actual level
$p^n$ is a failure at manuscript index $\nu=n-1$, so requires
$\nu_{\min}\ge n$. The largest $n\ge2$ with
$p^{2n-3}<e$ is the displayed integer, whenever that integer is
at least two. In the remaining range it equals one, and the
existing level-one height-two examples of
Proposition~\ref{prop:B-RD18}, with ramification a multiple of
$p-1$, give the same lower bound. Adding equal good height-one
factors proves every $h\ge2$ without changing either comparison.
\end{proof}

\begin{remark}[What the threshold does and does not settle]
\label{rem:B-RD63-scope}
\status{Candidate scope statement.}
The necessary bound assumes $p\nmid e$ and compares a whole
potentially supersingular group to a filtered height-one group.
It is not an upper bound for comparisons between two potentially
supersingular groups, and does not determine the unrestricted
least B.4 index outside the earlier exact cases. The existence
construction requires $p-1\mid e$ but not $p\nmid e$; it removes
the upper restriction $e<p^{2n-2}$ from B-RD62. The finite-character
and filtration conventions are exactly those of that theorem.

The exact supplement
\path{code/reduction_descent/verify_cyclic_torsion_threshold.py}
checks $433130$ Newton polygons and $1299390$ rational chain steps
for six primes with $a\le400$, $943$ explicit arcs, $5706$ weight
inequalities, and $14360$ integer-index comparisons. It passed in
$1.3598105860000942$ seconds at 20:00:45 UTC. The finite ranges
and actual adjusted choices of $A$ are retained in the output;
no infinite-field or BT existence assertion is inferred from them.
\end{remark}
